\subsubsection*{4.5.}
设 $f:X^*\longrightarrow D^*$ 是 $D^*$ 上的拓扑空间，$F$ 是 $X^*$ 上的 $f^\circ\mathcal O_{D^*}$-模层，$\bar F$ 是它在
\[
\bar X^*=X^*\times_{D^*}\widetilde D^*.
\]
上的逆像。

设 $z:D\hookrightarrow\mathbb C$ 是全纯嵌入，$z(0)=0$。下文总允许缩小 $D$；对足够小的圆盘，$z$ 的存在性是显然的。在 $\widetilde D^*$ 上选择函数 $\log z$ 使 $\exp(\log z)=z$，并令
\[
z^\alpha=\exp(\alpha\log z)
\qquad(\alpha\in\mathbb C).
\]

群 $\pi_1(D^*)$ 通过搬运结构作用于 $H^0(\bar X^*,\bar F)$。若一个元素 $f\in H^0(\bar X^*,\bar F)$ 在 $\pi_1(D^*)$ 作用下的各变换张成一个有限维复向量空间 $V(f)$，则称其为\underline{有限决定型}（finite determination）。设 $T$ 是 $\pi_1(D^*)$ 的正生成元。截面 $f$ 为有限决定型，当且仅当存在非零多项式 $P\in\mathbb C[X]$ 使 $P(T)(f)=0$。此时具有这一性质的最低次数多项式，既是 $T$ 在 $V(f)$ 上作用的最小多项式，也是其特征多项式。

若存在 $A$ 使
\[
(T-1)^A f=0.
\]
则称 $f$ 为\underline{幺幂有限决定型}。在这些截面构成的空间上定义
\[
N=\frac{\log T}{2\pi i}
\]
其作用由有限和
\[
N(f)=-\frac1{2\pi i}
\sum_{n>0}\frac{(1-T)^n}{n}\,f.
\]
给出。
若 $f$（分别为 $g$）是 $H^0(D^*,\mathcal O)$（分别为 $H^0(\bar X^*,\bar F)$）中的幺幂有限决定型截面，则
\[
N(fg)=N(f)g+fN(g),
\]
因为自同构的对数是导子。

若存在 $A$ 和 $B>0$ 使
\[
(T^B-1)^A f=0.
\]
则称 $f$ 为\underline{拟幺幂有限决定型}。

\subsubsection*{例 4.6.}
对 $X=D^*$ 和 $F=\mathcal O$，有
\[
\bigl(T-\exp(2\pi i\alpha)\bigr)^{k+1}
  \bigl(z^\alpha(\log z)^k\bigr)=0,
\]
\[
N(\log z)=1,
\]
以及
\[
N\left(\frac{(\log z)^k}{k!}\right)
=\frac{(\log z)^{k-1}}{(k-1)!},
\qquad
(0\text{ 当 }k=0).
\]

\subsubsection*{命题 4.7.}
\underline{截面 $f\in H^0(\bar X^*,\bar F)$ 为有限决定型，当且仅当它可写成}
\begin{equation}\tag{4.7.1}
f=\sum_{\alpha,k}f_{\alpha,k}\,
z^\alpha\frac{(\log z)^k}{k!},
\end{equation}
\underline{其中 $0\leq\operatorname{Re}(\alpha)<1$，$k\in\mathbb N$，且各 $f_{\alpha,k}\in H^0(X^*,F)$ 几乎全为零。对有限决定型的 $f$，分解 (4.7.1) 唯一。}

由 4.6，每个形如 (4.7.1) 的截面都为有限决定型。反之，设 $f$ 为有限决定型，且 $P(T)(f)=0$。写成
\[
P(X)=\prod_{\beta\in\mathbb C}(X-\beta)^{n'(\beta)}
=\prod_{0\leq\operatorname{Re}(\alpha)<1}
 \bigl(X-\exp(2\pi i\alpha)\bigr)^{n(\alpha)},
\]
其中 $n(\alpha)=n'(\exp(2\pi i\alpha))$ 几乎全为零。将 $\mathbb C[X]$-模 $V(f)$ 分解为其同型分量，得到唯一的分解
\[
f=\sum_\alpha f'_\alpha
\]
使
\[
\bigl(T-\exp(2\pi i\alpha)\bigr)^{n(\alpha)}f'_\alpha=0.
\]

选择多项式 $Q_\alpha\in\mathbb C[X]$，使其模 $\bigl(X-\exp(2\pi i\alpha)\bigr)^{n(\alpha)}$ 同余于 $1$，且对每个 $\alpha'\ne\alpha$，模 $\bigl(X-\exp(2\pi i\alpha')\bigr)^{n(\alpha')}$ 同余于 $0$。于是 $f'_\alpha=Q_\alpha(T)(f)$。令
\begin{equation}\tag{4.7.2}
f_\alpha=z^{-\alpha}f'_\alpha
=z^{-\alpha}Q_\alpha(T)(f).
\end{equation}
由于
\[
T(z^{-\alpha}f'_\alpha)
=\exp(-2\pi i\alpha)z^{-\alpha}T(f'_\alpha),
\]
有
\[
(T-1)^{n(\alpha)}f_\alpha=0.
\]

若
\[
f_\alpha=\sum_k f_{\alpha,k}\frac{(\log z)^k}{k!},
\qquad f_{\alpha,k}\in H^0(X^*,F),
\]
则由 4.6 及恒等式 $\sum_j(-1)^j\binom{k}{j}=\delta_{0,k}$ 得
\begin{equation}\tag{4.7.3}
f_{\alpha,k}
=\sum_{j\geq0}(-1)^j\frac{(\log z)^j}{j!}
  N^{j+k}f_\alpha,
\end{equation}
这证明唯一性。

反之，若 $f_\alpha$ 是幺幂有限决定型截面，则 (4.7.3) 定义的各 $f_{\alpha,k}$ 在 $k$ 足够大时为零（因为 $k$ 足够大时 $N^k f_\alpha=0$），且在 $T$ 下不变；因而它们可等同于 $X^*$ 上 $F$ 的截面。等价地，须证明 $Nf_{\alpha,k}=0$。将恒等式
\[
\begin{aligned}
N\sum_{j\geq0}(-1)^j\frac{(\log z)^j}{j!}N^j
&=\sum_{j>0}(-1)^j\frac{(\log z)^{j-1}}{(j-1)!}N^j\\
&\quad+\sum_{j\geq0}(-1)^j\frac{(\log z)^j}{j!}N^{j+1}=0.
\end{aligned}
\]
应用于 $N^k f_\alpha$。最后，
\[
\sum_{k\geq0}f_{\alpha,k}\frac{(\log z)^k}{k!}
=\sum_{j,k\geq0}(-1)^j
 \frac{(\log z)^{j+k}}{j!k!}N^{j+k}f_\alpha
\]
\[
=\sum_{n\geq0}
 \left(\sum_j(-1)^j\binom nj\right)
 \frac{(\log z)^n}{n!}N^n f_\alpha
=f_\alpha,
\]
这就证明了 4.7。

\subsubsection*{注 4.8.1.}
4.7 中的条件 $\operatorname{Re}(\alpha)\in[0,1)$ 可换为 $\alpha\in C$，其中 $C$ 是 $\mathbb C\longrightarrow\mathbb C/\mathbb Z$ 的任意基本域。

\subsubsection*{注 4.8.2.}
设 $f$ 为有限决定型，其分解为 (4.7.1)。对 $0\leq\operatorname{Re}(\alpha)<1$ 及满足某个 $f_{\alpha,k}\ne0$ 的 $\alpha$，令
\[
m(\alpha)=1+\sup\{k\mid f_{\alpha,k}\ne0\}.
\]
若所有 $f_{\alpha,k}$ 均为零，则令 $m(\alpha)=0$。$T$ 在 $V(f)$ 上作用的特征多项式为
\[
\operatorname{Pc}(T,V(f))
=\prod_\alpha\bigl(X-\exp(2\pi i\alpha)\bigr)^{m(\alpha)}.
\]
特别地，$f$ 为拟幺幂有限决定型，当且仅当 (4.7.1) 中的各 $\alpha$ 均可取在 $\mathbb Q$ 中。